% capitulos/cap5/isr.tex -- §5.7: renderização, cache e ISR (revalidação, visibilidade).
% Parte do Capítulo 5 (capitulos/cap5_arquitetura.tex); os rótulos \label e os números de seção não mudaram.

\section{Renderização, Cache e ISR}\label{sec:arq-isr}

\subsection{Estratégia por Rota}

A Tabela~\ref{tab:rotas-render} define o modo de renderização de cada conjunto de rotas. O projeto usa a ISR \emph{sem} Cache Components (o recurso \arq{cacheComponents} não está habilitado em \arq{next.config.mjs}), conforme o guia correspondente da documentação do Next.js 16 \cite{nextjsdocs}.

\begin{table}[htb]
\caption{Modo de renderização por conjunto de rotas}\label{tab:rotas-render}
\small
\begin{tabularx}{\linewidth}{@{}L{3.3cm}L{3.0cm}Y@{}}
\toprule
\textbf{Rota} & \textbf{Modo} & \textbf{Regras}\\
\midrule
\arq{/} (portal) & Estático com ISR & Lê o texto do portal e o diretório (RN28); invalidado só pela configuração do site e por mudanças de espaço que afetam o diretório (\texttt{revalidatePath("/")}, literal)\\
\addlinespace
\arq{/[handle]} e \arq{/[handle]/[...slug]} (M17) & Estático com ISR & \arq{generateStaticParams} lista só os caminhos \emph{publicados} de espaços \emph{ativos} (RN15, RN23), sem consultar o banco no \en{build} (\adr{ADR-021}); \texttt{revalidate} alto (3600 s) como rede de segurança; novos caminhos são gerados sob demanda; revalidação sob demanda por caminhos literais (\adr{ADR-030})\\
\addlinespace
\arq{/[...slug]} (hoje, até o M17) & Estático com ISR & Rota atual, sem espaço; substituída pelas duas anteriores\\
\addlinespace
\arq{/login} & Estático ou dinâmico & Formulário; sem dados sensíveis\\
\addlinespace
\arq{/admin/**} & Dinâmico & Lê a sessão por \en{cookies}; nunca em cache\\
\addlinespace
\arq{/api/**} & Dinâmico & Route Handlers; respostas sem cache\\
\bottomrule
\end{tabularx}
\fonte{Elaborado pelos autores.}
\end{table}

\subsection{Regras de Revalidação}

A Tabela~\ref{tab:revalidacao} implementa as regras RN10 e RN16 \emph{por espaço} (\adr{ADR-030}; módulo M17). A invalidação é feita em Route Handlers com \arq{revalidatePath}, função que só pode ser chamada no servidor, não em Componentes Cliente nem no Proxy. Em Route Handlers, a função marca o caminho para revalidação, que ocorre na próxima visita \cite{nextjsdocs}.

\textbf{Por que caminhos literais.} Medido no Next.js 16.3.7 (\texttt{next build} e \texttt{next start}, dois espaços, cabeçalho \texttt{x-nextjs-cache}): \texttt{revalidatePath('/a', 'layout')} \emph{não invalida nada} em um segmento dinâmico; \texttt{revalidatePath('/[handle]', 'layout')} invalida \emph{todos} os espaços; \texttt{revalidatePath('/', 'layout')} invalida tudo; só o caminho literal de página (\texttt{revalidatePath('/a/x')}) funciona, e é barato (2000 chamadas em 58~ms). Por isso o servidor calcula os caminhos públicos do nó e dos descendentes \emph{antes e depois} da escrita, a partir da árvore do dono (\texttt{getPublicPaths}, do M15), porque renomear ou mover deixa em \en{cache} o HTML antigo. Com até 500 conteúdos por dono (RN25), são cerca de 500 chamadas por evento, abaixo de 10~ms. \textbf{Nunca} se usa \texttt{revalidatePath('/', 'layout')} por mudança de um espaço ou do site. O comportamento depende da versão do \en{framework} e deve ser conferido a cada atualização (CT44).

\textbf{Decisões do M-015 (05/10/2026).} (i) Os caminhos das coleções \emph{ancestrais} também são invalidados, porque uma coleção redireciona ao primeiro documento publicado (RN14) e criar, despublicar, mover ou excluir esse documento muda a página da coleção. (ii) Os eventos \emph{estruturais} invalidam o \textbf{espaço inteiro do dono}, por caminhos literais (\texttt{getPublicPaths} da raiz do espaço, no máximo 500 conteúdos, RN25), e não só o nó e os descendentes: a \en{sidebar} de cada página mostra a árvore inteira, e as páginas irmãs ficariam velhas até o \texttt{revalidate} de 3600~s (custo desprezível, \adr{ADR-030}). Só a edição do corpo invalida apenas a própria página. \emph{Implementado no merge M-019 (05/10/2026):} \texttt{revalidateStructure} captura o espaço inteiro antes e depois da escrita e invalida a união (que já traz as coleções ancestrais); só o corpo invalida só a página. O cabeçalho \texttt{x-nextjs-cache} do CT44 só se mede com as páginas públicas (Parte 3b). (iii) O portal é invalidado também quando muda o diretório (RN28): raiz criada, excluída, publicada, despublicada ou movida, \en{handle}, suspensão e nome do usuário. \textbf{O caminho é \texttt{revalidatePath("/")}, e não \texttt{revalidatePath('/', 'page')}} (medido no merge M-021, sobre o \texttt{next start}): com o portal em um grupo de rotas (\texttt{(portal)}, que dá o layout só a \texttt{/}), a página se chama \texttt{/(portal)/page}, e o tipo \texttt{page} de \texttt{/} não a alcança (o \texttt{x-nextjs-cache} seguia \texttt{HIT} depois do \texttt{PUT}); o caminho literal sem tipo a regenera (\texttt{MISS}) e continua sendo literal: o tipo \texttt{layout} segue proibido.

\begin{table}[htb]
\caption{Eventos de mutação e a revalidação correspondente}\label{tab:revalidacao}
\small
\begin{tabularx}{\linewidth}{@{}L{5.2cm}Y@{}}
\toprule
\textbf{Evento} & \textbf{Ação de revalidação}\\
\midrule
Criar ou excluir conteúdo; alterar título, \en{slug}, pai, ordem ou tipo; mover/reordenar & \textbf{Todo o espaço do dono}: página inicial e todos os caminhos públicos, \emph{antigos e novos} (a \en{sidebar} de cada página lista a árvore inteira, então as páginas irmãs também ficam velhas). Se o evento é sobre uma raiz (criada, excluída, publicada, despublicada ou movida de ou para a raiz), o portal também (RN28)\\
\addlinespace
Alterar o estado de publicação (\texttt{DRAFT} $\leftrightarrow$ \texttt{PUBLISHED}) de qualquer nó, inclusive por moderação & O mesmo, com os caminhos \emph{antigos e novos}: a visibilidade do nó e de seus descendentes altera a navegação e o conjunto de páginas (RN16)\\
\addlinespace
Alterar somente o corpo de um documento (estado inalterado) & Só \texttt{revalidatePath(caminho do documento)}, e só se a página é pública\\
\addlinespace
Alterar a aparência do dono, o \en{handle} (com o antigo e o novo), a suspensão ou a exclusão da conta (e, na transferência, o espaço do destino) & Todo o espaço: página inicial e todos os caminhos públicos do dono; \en{handle}, suspensão, exclusão e o nome do usuário também invalidam o portal\\
\addlinespace
Alterar a configuração do site ou o portal (nome, logotipo, cores, links, texto, diretório) & \texttt{revalidatePath("/")} (o caminho literal do portal; ver a nota do M-021 acima). \emph{Integrado (épico v1, \adr{ADR-038}, item 6; M-041):} se a mudança é de uma unidade que os usuários herdam (a paleta clara, a paleta noturna ou o nome do sistema), também todo o espaço de cada usuário que herda aquela unidade; se muda o par de ícones do botão de modo (\texttt{modeIconPair}), o espaço de \emph{todos} (o botão aparece em todo espaço e o par é só do sistema). \emph{Integrado (M-042):} trocar ou remover o logotipo ou o \en{favicon} do sistema revalida o portal e os espaços que herdam aquela imagem (mesma operação, unidades \texttt{logo} e \texttt{favicon})\\
\addlinespace
\emph{Integrado (épico v1, \adr{ADR-034}; M-042):} enviar, trocar ou remover o logotipo ou o \en{favicon} de um dono & Todo o espaço do dono (o \en{favicon} sai no \texttt{generateMetadata} e o logotipo, na barra superior de cada página); se for do site, o portal e os espaços que herdam a imagem\\
\emph{Alvo (v1.1, \adr{ADR-040}):} salvar, despublicar ou excluir a página inicial de um espaço, ou ligar e desligar a lista de itens & Só a página inicial do espaço (\texttt{revalidatePath} de \texttt{/\{handle\}}); no sistema, o portal \texttt{/}. A folha do KaTeX só é incluída quando o texto tem fórmula (RNF02)\\
\addlinespace
\emph{Integrado (M-042):} enviar, trocar ou remover o avatar; a moderação do avatar & O espaço do dono e o portal (o avatar aparece no diretório)\\
\addlinespace
\emph{Alvo (\adr{ADR-035}):} substituir o rodapé de um dono & Todo o espaço do dono (o rodapé está em todas as páginas dele); o rodapé do site revalida o portal\\
\bottomrule
\end{tabularx}
\fonte{Elaborado pelos autores, com base em RN10, RN16 e na \adr{ADR-030}.}
\end{table}

\textbf{O layout continua estático (épico v1).} Nenhuma das mudanças do épico torna uma página dinâmica: o modo de cor está no \en{cookie} \texttt{vitrine-modo}, lido só pelo \en{script} do \texttt{<head>} (nunca pelo servidor; \adr{ADR-036}); o item ``Administração'' das páginas públicas depende do \en{cookie}-marcador \texttt{vitrine-logado}, lido no navegador (\adr{ADR-037}); a mescla da aparência é feita na geração (\adr{ADR-038}); e \texttt{GET /media/<id>} é um \en{Route Handler} fora das páginas, com resposta imutável. O ensaio HTTP do Arquiteto confere \texttt{x-nextjs-cache} das páginas públicas depois do épico (CT47 e CT64).

\textbf{Revalidar os espaços que herdam o sistema (integrado, M-041).} Só o administrador muda o sistema, e raramente; a rota compara a configuração antes e depois (\texttt{changedSystemUnits}) e \texttt{revalidateInheritingSpaces} invalida, por caminhos literais (\adr{ADR-030}), o espaço inteiro de cada usuário \emph{ativo} que herda ao menos uma das unidades que mudaram (\texttt{inherits} de \texttt{getEffectiveAppearance}); quem definiu a própria paleta de um modo não é tocado por uma mudança daquele modo. \emph{Custo medido} (\texttt{MEDICAO=1}, \texttt{revalidatePath} simulado: mede o banco e a montagem dos caminhos, e não a invalidação do Next.js): 20 espaços de 10 páginas, 110 URLs em 76~ms; 100 espaços de 20 páginas (metade herda), 1050 URLs em 234~ms; 200 de 50 (metade herda), 5100 URLs em 394~ms; todos os 200 (o caso do \texttt{modeIconPair}), 10200 URLs em 245~ms. A operação é feita dentro da requisição do administrador e não falha a gravação, que já foi confirmada (RNF10). \emph{Limite conhecido, aceito pelo Engenheiro Chefe em 09/10/2026:} só entram os espaços que \texttt{listPublicSpaces} devolve (ao menos uma raiz publicada, RN28); o início de um espaço \emph{sem} raiz publicada, se tiver sido visitado, fica em cache e só se atualiza pelo \texttt{revalidate} de 3600~s (a correção é uma função do Dados que liste os \emph{handles} ativos, no \en{backlog} da v1.1; o risco é baixo, porque esse início quase não tem conteúdo a mostrar além da moldura).

\textbf{Latência observada da revalidação: era o limite de escritas, não o \en{cache}.} Nos primeiros ensaios do Arquiteto sobre o \en{build} de produção (Next.js 16.3.7, \texttt{next start}, uma instância), a leitura \emph{seguinte} a um despublicar ou republicar às vezes ainda servia o estado antigo, com \texttt{x-nextjs-cache} \texttt{HIT}, por 2 a 20~s. O Core provou (M-042, P-019) que a causa não é a revalidação: a partir da 61ª escrita no minuto, o \texttt{PATCH} responde \textbf{429} (60 escritas por minuto por usuário, RN25) e \emph{nada muda no servidor}, de modo que a página fica como estava; os ensaios do Arquiteto não conferiam o status da escrita. \emph{Confirmação independente do Arquiteto} (\en{build} de produção no Linux, 45 ciclos de despublicar e republicar, 90 escritas): as 67 escritas aceitas (200) refletiram na \emph{primeira} leitura seguinte, sem exceção; as 23 recusadas (429, todas depois da 60ª) não alteraram a página. A revalidação é, portanto, concluída em segundos para toda escrita aceita (a ressalva de milissegundos está logo abaixo e fixa a redação da RN15), e o ensaio HTTP agora confere o status da escrita e espera o \texttt{Retry-After} (\texttt{write}). A mesma causa explicava a oscilação do \texttt{e2e/usuarios.spec.js}: o \texttt{DELETE} e o \texttt{/ban} do grupo \texttt{users} (20 por minuto) voltavam 429, gastos pelos \en{specs} anteriores, e a interface não esperava o \texttt{Retry-After}; o Core a corrigiu com um segundo administrador só para as sessões de navegador.

\textbf{Ressalva de milissegundos (merge M-043).} A confirmação acima vale para os segundos, não para o milissegundo: em 3 de 5 execuções do ensaio completo (as 5 primeiras, com a verificação estrita da primeira leitura), uma leitura feita logo depois de um \texttt{PATCH} \emph{aceito} (200) ainda recebeu a versão antiga, com \texttt{x-nextjs-cache} \texttt{HIT}, nos dois sentidos (despublicar, republicar e reativar um espaço suspenso), e a leitura seguinte, \textbf{33 e 249~ms depois}, já era a nova. Não é o limite de escritas (o status era 200) e não repetiu em 40 ciclos isolados de leitura, escrita e leitura (80 escritas), nem nas 5 execuções seguintes, em que o ensaio passou a esperar o estado e a registrar o atraso; a causa não foi isolada (hipótese: a regeneração da página, disparada pela leitura imediatamente anterior, termina depois da invalidação). \textbf{Redação adotada (decisão do Engenheiro Chefe, 09/10/2026; RN15 e UC):} a requisição \emph{imediatamente seguinte} a uma escrita aceita pode, por milissegundos, ainda trazer a versão anterior da página, e \emph{a partir da requisição seguinte a essa} vale a nova; ``a primeira visita depois da escrita é sempre a nova'' \emph{não} é garantido. A frase antiga, ``a partir da próxima requisição'', foi substituída no Capítulo~\ref{cap:requisitos} (RN15 e caso de uso UC05), no CT21, no CT32 e no critério de aceite do M10. O ensaio HTTP trata a espera como informação (linha \texttt{INFO}), e não como falha, e registra o atraso observado. A atualização do \texttt{next} para a 16.4.0, recomendada pela auditoria de dependências (Seção~\ref{sec:arq-npm-audit}; correções de \emph{cache poisoning} de ISR), foi autorizada pelo Engenheiro Chefe em 09/10/2026 (pedido P-021 ao Core) e integrada no merge M-044. \textbf{Remedição sobre o código final (10/10/2026, Next.js 16.4.0, \en{build} de produção no Linux):} em 5 execuções completas do ensaio (496 verificações cada) a primeira leitura depois de toda escrita aceita já foi a nova (\texttt{INFO: nenhuma primeira leitura velha}), e em 40 ciclos isolados de leitura, escrita e leitura (80 escritas) também não houve leitura velha. Isso \emph{não} permite creditar a correção à 16.4.0: sob a 16.3.7 as mesmas condições da última série (5 execuções com o ensaio esperando o estado) também deram zero, e as 3 ocorrências em 5 vieram só das primeiras execuções; o fenômeno é intermitente, a causa segue não isolada e a redação da RN15 (``a requisição imediatamente seguinte pode, por milissegundos, ainda trazer a versão anterior'') permanece, porque uma ausência de ocorrências não prova a ausência da janela.

\textbf{Páginas inexistentes.} Uma página inexistente fica em \en{cache} (404 com \texttt{HIT}) e ocupa cerca de 12 a 16~kB e 7 arquivos por URL, sem expurgo (200 URLs aleatórias acrescentaram 1400 arquivos). Para o primeiro segmento que não é \en{handle} nem rota reservada, o Proxy faz \texttt{rewrite} para uma página dinâmica que responde \texttt{notFound()}: nenhum arquivo novo, e um \en{handle} criado em seguida funciona sem revalidar (CT45). O documento inexistente \emph{dentro de um espaço válido} não é coberto, e o risco, medido, é aceito com a mitigação operacional do Apêndice~\ref{ap:implantacao} (limite de taxa no \en{proxy} reverso e monitoramento do disco); O Proxy (M-024) reescreve ao 404 dinâmico (\arq{/_nao-encontrado}, segmento \texttt{\%5F}: nem \en{handle} nem \en{slug}, então sem novo nome reservado) o primeiro segmento que não é um \en{handle} \emph{ativo} e todo segmento fora da forma de \en{handle} ou de \en{slug} (maiúsculas, ponto, codificado, longo): nenhum arquivo em \arq{.next/} (200 URLs aleatórias, a mesma contagem de arquivos), um \en{handle} criado em seguida responde 200 na primeira visita, o espaço suspenso ou removido some antes do cache, e o 404 em maiúsculas, que em sistema de arquivos sem distinção de caixa (o Windows) colidia com a página real, deixa de envenená-la. Custo medido do Proxy (mediana e p95, dono com 500 conteúdos): \en{handle} ativo, 0,59~ms e 1,07~ms; \en{handle} mais documento, 4,30~ms e 8,02~ms, cerca de sete vezes; \en{handle} desconhecido, 0,42~ms; forma inválida, 0,04~ms. A consulta do documento no Proxy foi \emph{rejeitada}: custaria isso em toda visita para evitar só o 404 em \en{cache} de documento inexistente dentro de um espaço válido, risco aceito com a mitigação do Apêndice~\ref{ap:implantacao}. Uma falha do banco no Proxy deixa a requisição seguir (a página decide).

\subsection{Visibilidade por Estado de Publicação}\label{sec:arq-visibilidade}

A regra RN15 --- um conteúdo só é público se ele e todos os seus ancestrais estiverem \texttt{PUBLISHED} \emph{e o espaço do dono estiver ativo} (não suspenso, com \en{handle}; RN23) --- é aplicada em um único lugar, a camada de dados, e nunca nos componentes. Os pontos de aplicação são: (i) a árvore usada pela \en{sidebar}, pela barra superior e por \arq{generateStaticParams}, podada conforme o Código~\ref{lst:algoritmo-arvore}; (ii) a resolução do caminho nos Server Components, que responde 404 (\arq{notFound()}) para um nó ausente da árvore podada, de modo que um rascunho não é acessível nem por URL direta; e (iii) o mapa do site e os metadados (RF25). As rotas \arq{/api/contents} e o painel recebem a árvore completa \emph{do dono}, com o campo \texttt{status} de cada nó (RF26); o administrador vê a de outro usuário só no painel, para a moderação (RN19). Como o estado entra na invalidação do espaço (Tabela~\ref{tab:revalidacao}), a mudança de estado reflete-se nas páginas do dono a partir da segunda requisição depois da escrita (a primeira pode, por milissegundos, ainda trazer o estado anterior; Seção~\ref{sec:arq-isr}). Como o padrão do estado é \texttt{PUBLISHED} (\decisao{D-15}), a proteção contra a publicação acidental depende de a interface exibir o estado de forma explícita e de a API aceitar a criação já como rascunho (risco R11).

\vitrinefigura[tipo=diagrama, arquivo=seq-leitura-isr, agente=Arquiteto, largura=0.95,
  origem={doc/diagrams/src/seq-leitura-isr.puml},
  descricao={Sequência de leitura de uma página com os três caminhos do cache: HIT, STALE e MISS, incluindo a regeneração em segundo plano e o cabeçalho x-nextjs-cache.}]
  {Sequência de leitura de uma página com ISR}{fig:seq-leitura-isr}

\vitrinefigura[tipo=diagrama, arquivo=seq-publicacao-conteudo, agente=Arquiteto, largura=0.95,
  origem={doc/diagrams/src/seq-publicacao-conteudo.puml},
  descricao={Sequência de edição e publicação: editor, PATCH na API, validação de sessão e de entrada, atualização no banco e revalidação do caminho ou do layout.}]
  {Sequência de edição e publicação de um documento}{fig:seq-publicacao-conteudo}

\subsection{Limitações e Observabilidade}

A documentação do Next.js 16 registra ressalvas que condicionam a implantação: a ISR exige o \en{runtime} Node.js e não funciona em exportação estática; com mais de uma instância, o cache padrão em sistema de arquivos é por instância e a invalidação só atinge a instância que recebeu a chamada; e o Proxy não é executado em requisições de ISR sob demanda \cite{nextjsdocs}. Por isso, a implantação usa \textbf{uma única instância}, isto é, um único contêiner (o que também convém ao SQLite). Em testes, o cabeçalho \arq{x-nextjs-cache} (\texttt{HIT}, \texttt{STALE}, \texttt{MISS}, \texttt{REVALIDATED}) permite verificar o comportamento do cache (Capítulo~\ref{cap:testes}).
